% Artifact Evaluation (AE) appendix — draft, NOT for initial paper submission.
% Per SC26 guidelines (inherited from SC25), the AE appendix is submitted
% only after paper acceptance when requesting reproducibility badges.
% To assemble the full AD+AE document, paste the body of this file
% (everything between \appendixAE and \end{document}) at the end of
% artifact_description.tex, immediately before its \end{document}.

\documentclass[nonacm,sigconf]{acmart}

\usepackage{cite}
\usepackage{amsmath,amsfonts}
\usepackage{algorithmic}
\usepackage{graphicx}
\usepackage{textcomp}
\usepackage{xcolor}
\usepackage{booktabs}
\usepackage{tagging}
\usepackage{url}
\usepackage{listings}

\lstset{
  basicstyle=\footnotesize\ttfamily,
  breaklines=true,
  frame=single,
  columns=fullflexible,
}

\usepackage{sc26repro}
\pagenumbering{gobble}

\begin{document}

\twocolumn[%
{\begin{center}
\Huge
Appendix: Artifact Evaluation (draft)
\end{center}}
]

%%%%%%%%%%%%%%%%%%%%%%%%%%%%%%%%%%%%%%%%%%%%%%%%%%%%%
%  AE Appendix
%%%%%%%%%%%%%%%%%%%%%%%%%%%%%%%%%%%%%%%%%%%%%%%%%%%%%
\appendixAE

% ---- A1: DALIA/ADELIA ----
\arteval{1}

\artin

$A_1$ (the ADELIA framework) is installed as part of the
benchmark suite ($A_3$) via a single command on CSCS Alps:

\begin{lstlisting}[language=bash]
git clone <artifact-repo> adelia-artifact
cd adelia-artifact
salloc -N1 --gpus-per-task=1 -A <account> -t 30 -p debug
uenv run --view=modules prgenv-gnu/24.11:v1 -- \
    bash scripts/install_alps.sh
\end{lstlisting}

\noindent
\texttt{install\_alps.sh} creates the \texttt{adelia-env} conda
environment, installs JAX 0.8.1 + CUDA plugin, CuPy 13.6, builds
\texttt{mpi4py} and \texttt{mpi4jax} from source against Cray MPICH
with CUDA, and installs $A_1$ and $A_2$ as editable packages.
See \texttt{INSTALL.md} for the manual step-by-step path and
\texttt{INSTALL.md\#Troubleshooting} for common pitfalls
(cuSolver/CuPy import order, \texttt{libdevice.10.bc}, MPI linking).

\artcomp

ADELIA exposes its gradient routines through the \texttt{DALIA}
Python API:
\begin{lstlisting}[language=bash]
conda activate adelia-env
python -c "
import jax, jax.numpy as jnp
jnp.linalg.cholesky(jnp.eye(2)).block_until_ready()  # GPU warmup
import os; os.environ['ARRAY_MODULE'] = 'cupy'
from dalia.core.jax_autodiff import \
    create_pure_jax_objective_distributed_coregional_splitjit
print('ADELIA AD entry point imported OK')
"
\end{lstlisting}

\noindent
The JAX-first, CuPy-second import order is mandatory (see $A_3$
execution for the full evaluation workflow).

\artout

The framework has no standalone outputs; its correctness and
performance are evaluated through $A_2$'s test suite and $A_3$'s
benchmark and validation jobs. See \texttt{DALIA/tests/README.md}
for framework-level unit tests.


% ---- A2: Serinv ----
\arteval{2}

\artin

Installed as part of $A_1$. Manual install:
\begin{lstlisting}[language=bash]
conda activate adelia-env
cd serinv && pip install -e .
\end{lstlisting}

\artcomp

Run the sequential and distributed test suites:
\begin{lstlisting}[language=bash]
cd serinv
pytest tests_algs/ tests_wrappers/
mpirun -n 4 pytest --with-mpi tests_wrappers/
\end{lstlisting}

\noindent
The sequential suite exercises factorization, solve, and selected
inversion against dense reference implementations. The distributed
suite verifies that nested-dissection partitioned results match
sequential execution on block-tridiagonal and BTA matrices.

\artout

All tests should pass ($\sim$10~min sequential, $\sim$5~min
distributed). A passing test run substantiates $C_1$ (structured
BTA algorithms producing numerically identical results to dense
solvers) and $C_3$ (distributed two-phase correctness).


% ---- A3: Benchmark suite ----
\arteval{3}

\artin

Included in the artifact repository. After completing the $A_1$
install step above:

\begin{lstlisting}[language=bash]
cd adelia-artifact
./scripts/download_data.sh            # ~1.5 GB into data/
# Edit two values in experiments/common/env_setup.sh:
#   SBATCH_ACCOUNT=<your CSCS account>
#   CONDA_ENV=adelia-env
\end{lstlisting}

\artcomp

The benchmark suite has three task groups. Reviewers with a
strict 8-hour budget should run only the \textbf{reviewer path}
below, which exercises every code path on a 1-GPU subset and
regenerates every paper figure from the reference CSVs shipped
in \texttt{plotting/data/}:

\begin{lstlisting}[language=bash]
# (1) Correctness: ~10 min, 1 GPU
sbatch validation/gst_small/validate.sbatch
sbatch validation/gst_coreg3_small/validate.sbatch

# (2) Functional: ~30 min, 1 GPU
sbatch experiments/fig4_wallclock/gst_small/run.sbatch
sbatch experiments/fig4_wallclock/gst_coreg3_small/run.sbatch

# (3) Figure regeneration: ~5 min, no GPU
cd plotting && bash generate_all_figures.sh
\end{lstlisting}

\noindent
Full reproduction ($\sim$28 GPU-hours, required for the Reproduced
badge) is documented per-experiment in the \texttt{README.md} files
under \texttt{experiments/tab2\_structure\_comparison/},
\texttt{fig2\_wallclock/}, \texttt{fig4\_framework\_decomposition/},
\texttt{fig5\_scaling\_study/}, \texttt{fig6\_resource\_efficiency/},
and \texttt{fig7\_performance\_analysis/}.
Task dependencies: $T_1$ (validation) is independent; $T_2$
(experiments) directories are independent of one another; $T_3$
(plotting) consumes either fresh $T_2$ CSVs or the reference CSVs
in \texttt{plotting/data/}.

\artout

Reviewer path outputs:
\begin{itemize}
  \item Each \texttt{validate.sbatch} writes a PASS/FAIL summary to
    \texttt{validation/<model>/outputs/*.out}. Pass criteria are
    documented in \texttt{validation/README.md}.
  \item \texttt{fig2\_wallclock/<model>/run.sbatch} writes
    per-gradient timing CSVs into each model's \texttt{outputs/}
    directory.
  \item \texttt{plotting/generate\_all\_figures.sh} writes the
    paper's PDF figures into \texttt{plotting/figures/}: Table~2,
    Figures~3 (convergence), 4 (wall-clock), 5 (scaling),
    6 (framework decomposition), 7 (resource/energy),
    8 (performance analysis).
\end{itemize}

\noindent
Expected correlations to contributions:
the validation PASS confirms $C_1$--$C_3$ numerically;
the wall-clock CSVs (and the regenerated Figure~4) demonstrate
$2.4$--$50.9\times$ per-gradient speedups over FD ($C_1$--$C_3$);
the regenerated Figures~5--7 substantiate $C_4$ (scaling,
framework decomposition, energy savings).

\end{document}
